\chapter{Muffins}

Muffins are a Kiwi classic — equally at home in lunchboxes, on café counters, or warm from the oven with a cuppa.  
They can be sweet or savoury, simple or indulgent, but they’re always quick to bake and versatile enough to suit every taste.  
This chapter brings together favourites from rich chocolate and tangy lemon to seasonal feijoa, alongside hearty savoury options like corn and spinach with feta.  
Perfect for breakfast, snacks, or a grab-and-go treat, there’s a muffin here for every occasion.  

\minitoc

\section*{Sweet Muffins}
\addcontentsline{toc}{section}{Sweet Muffins}

\begin{recipebox}{Chocolate Muffins}{12}{40 min}
\begin{ingredients}
  \item 2 eggs
  \item 1 cup White Sugar (225 g)
  \item 1 cup Standard Grade Flour (150 g)
  \item 6 Tbsp cocoa powder
  \item 2 tsp Baking Powder
  \item \nicefrac{1}{2} tsp salt
  \item \nicefrac{2}{3} cup Milk (160 mL)
  \item 165 g butter, melted
  \item 1 tsp vanilla essence
  \item \nicefrac{2}{3} cup chocolate chips (130 g), tossed with 1 Tbsp flour (optional)
\end{ingredients}

\begin{method}
  \item Preheat oven to 180\degree C bake. Grease or line 6 extra-large muffin tins or 12–15 medium muffin tins.
  \item Beat the eggs and sugar together until well combined.
  \item In a separate bowl, mix the flour, cocoa powder, baking powder, and salt. Add to egg mixture and stir to just combine.
  \item Add the milk, melted butter, and vanilla essence. Fold in gently until just combined. If using, fold in the flour-coated chocolate chips.
  \item Spoon mixture into prepared tins until approximately two-thirds full.
  \item Bake for 18–20 minutes (regular muffins) or 30 minutes (extra-large muffins), until muffins spring back when lightly pressed or a skewer inserted into the centre comes out clean.
  \item Leave to cool in tins for 10 minutes before transferring to a wire rack. Serve warm or at room temperature.
\end{method}

\begin{notes}
For extra indulgence, top with a sprinkle of chocolate chips before baking.  
This recipe can be made with dark, milk, or white chocolate chips.  
\end{notes}
\end{recipebox}

\begin{recipebox}{Banana Chocolate Chip Muffins}{12}{30 min}
\begin{ingredients}
  \item 1 \nicefrac{1}{2} cups all-purpose flour (180 g)
  \item 1 tsp baking soda
  \item \nicefrac{1}{2} tsp salt
  \item \nicefrac{2}{3} cup sugar (113 g)
  \item 1 cup semisweet chocolate chips (plus extra for topping, optional)
  \item 3 medium ripe bananas (about 1 cup mashed)
  \item \nicefrac{1}{3} cup melted butter, slightly cooled (75 g)
  \item 1 large egg
  \item 1 tsp vanilla essence (5 mL)
\end{ingredients}

\begin{method}
  \item Preheat oven to 190\degree C (375\degree F). Line a muffin tin with papers.  
  \item In a bowl, whisk together flour, baking soda, salt, and sugar. Set aside.  
  \item In a large bowl, mash the bananas. Add egg, melted butter, and vanilla. Mix to combine.  
  \item Add the dry mixture to the wet ingredients and fold gently until almost combined.  
  \item Fold in chocolate chips and mix until evenly distributed.  
  \item Spoon batter into muffin tin and top with extra chocolate chips if desired.  
  \item Bake for 17–20 minutes, or until a toothpick inserted in the centre comes out clean.  
\end{method}

\begin{notes}
For best results, use ripe bananas (skin mostly brown).  
Optional additions: ½ cup walnuts or pecans, or ½ tsp cinnamon.  
Muffins keep 5–7 days in an airtight container or can be frozen for up to 3 months.  
To ripen bananas quickly: bake whole bananas (in skins) at 180\degree C for 10 minutes, let cool, then mash.  
\end{notes}
\end{recipebox}



\begin{recipebox}{Lemon Muffins}{12}{30 min}
\begin{ingredients}
  \item 2 cups Self-Raising Flour (300 g)
  \item \nicefrac{3}{4} cup Caster Sugar (170 g)
  \item 1 egg
  \item 1 cup Milk (250 mL)
  \item Finely grated zest of 1 lemon
  \item 60 g butter, melted (or Olivani spread)
\end{ingredients}

\subsection*{Lemon Drizzle Topping}
\begin{itemize}[leftmargin=2em]
  \item \nicefrac{1}{2} cup lemon juice (80 mL)
  \item \nicefrac{1}{2} cup Caster Sugar (75 g)
  \item 60 g butter, melted
\end{itemize}

\begin{method}
  \item Preheat oven to 180\degree C fan bake (200\degree C regular bake). Grease or line a 12-hole muffin tin.
  \item Mix flour and caster sugar together in a large bowl.
  \item In a separate bowl, whisk egg, milk, lemon zest, and melted butter. Add to dry ingredients and stir until just combined.
  \item Pour mixture into prepared muffin tin.
  \item Bake for 15–20 minutes, until muffins spring back when lightly pressed.
  \item While muffins are baking, whisk together lemon juice, caster sugar, and melted butter until sugar dissolves.
  \item As soon as muffins are cooked, spoon lemon drizzle topping over them. Leave in the tin for 3 minutes, then transfer to a wire rack to cool.
\end{method}

\begin{notes}
These muffins are best enjoyed warm with extra lemon drizzle.  
For extra zing, add a teaspoon of lemon juice directly into the batter.  
\end{notes}
\end{recipebox}

\begin{recipebox}{Blueberry Muffins}{12}{35 min}
\begin{ingredients}
  \item 1 \nicefrac{1}{2} cups Milk (375 mL)
  \item 1 cup Caster Sugar (225 g)
  \item \nicefrac{1}{4} cup Olive Oil (60 mL)
  \item 1 egg
  \item 1 tsp lemon zest
  \item 2 tsp lemon juice
  \item 2 cups Standard Grade Flour (300 g)
  \item 4 tsp Baking Powder
  \item 1 \nicefrac{1}{2} cups blueberries, fresh or frozen (240 g)
  \item Icing Sugar, for dusting (optional)
\end{ingredients}

\begin{method}
  \item Preheat oven to 200\degree C bake (180\degree C fan bake). Line a 12-hole muffin tin with paper cases or grease well.
  \item In a large bowl, whisk together milk, sugar, oil, egg, lemon zest, and lemon juice.
  \item Sift flour and baking powder over the wet mixture. Fold in gently until just combined.
  \item Gently fold in blueberries, being careful not to over-mix.
  \item Divide mixture evenly between muffin cases.
  \item Bake for 17–20 minutes, until muffins are golden and spring back when lightly pressed.
  \item Transfer to a wire rack to cool. Dust with icing sugar before serving, if desired.
\end{method}

\begin{notes}
Fresh or frozen blueberries both work well — no need to thaw if frozen.  
Over-mixing will make the muffins tough and affect their rise.  
These are delicious served warm with a dusting of icing sugar.  
\end{notes}
\end{recipebox}

\begin{recipebox}{Feijoa Muffins}{12}{30 min}
\begin{ingredients}
  \item 1 cup chopped fresh feijoas (275 g)
  \item 75 g butter, melted
  \item 2 eggs
  \item \nicefrac{3}{4} cup White Sugar (170 g)
  \item 2 cups Standard Grade Flour (300 g)
  \item 2 tsp Baking Powder
  \item \nicefrac{1}{2} cup Meadow Fresh Milk (125 mL)
\end{ingredients}

\subsection*{Topping}
\begin{itemize}[leftmargin=2em]
  \item 2 Tbsp Demerara Sugar
  \item 1\nicefrac{1}{2} tsp cinnamon
\end{itemize}

\begin{method}
  \item Preheat oven to 200\degree C bake (180\degree C fan bake). Line a 12-hole muffin tin with paper cases or grease well.
  \item Cut feijoas in half, scoop out the flesh, and roughly chop until you have 1 cup (275 g).
  \item In a large bowl, melt butter. Stir in feijoas, eggs, and sugar.
  \item Add flour, baking powder, and milk. Fold together gently until just combined.
  \item Divide mixture evenly between muffin tin holes. Mix topping ingredients and sprinkle over muffins.
  \item Bake for 15–20 minutes, until muffins are golden and spring back when lightly pressed. Cool on a wire rack.
\end{method}

\begin{notes}
These muffins are a perfect way to use seasonal feijoas in autumn.  
Delicious served warm with butter, or packed into lunchboxes.  
You can freeze them for up to 3 months.  
\end{notes}
\end{recipebox}

\begin{recipebox}{Raspberry and White Chocolate Muffins}{12}{30 min}
\begin{ingredients}
  \item 2 \nicefrac{1}{2} cups Self-Raising Flour
  \item 100 g butter, chopped (or Olivani spread)
  \item 1 cup Caster Sugar
  \item 1 cup Milk
  \item 1 egg, lightly beaten
  \item 1 cup fresh or frozen raspberries
  \item \nicefrac{1}{2} cup white chocolate bites
  \item Icing Sugar, to dust (optional)
\end{ingredients}

\begin{method}
  \item Preheat oven to 190\degree C. Lightly grease or spray a 12-hole muffin tin.
  \item Place flour in a large bowl. Using your fingertips, rub butter into flour until well combined.
  \item Add sugar, milk, egg, raspberries, and white chocolate. Mix gently until just moistened.
  \item Spoon mixture evenly into muffin cups.
  \item Bake for 20 minutes, until golden. Leave to stand for 5 minutes before turning out onto a wire rack to cool.
  \item Dust with icing sugar before serving, if desired.
\end{method}

\begin{notes}
Frozen raspberries work just as well as fresh — add them straight from the freezer to avoid bleeding too much colour.  
Delicious served warm, with the white chocolate still slightly melted.  
\end{notes}
\end{recipebox}

\begin{recipebox}{Crunchy Lemon Poppyseed Muffins}{12}{25 min}
\begin{ingredients}
  \item 2 cups self-raising flour
  \item 1 cup White Sugar
  \item Finely grated rind of 2 lemons
  \item 100 g butter (or Olivani spread), melted
  \item 2 large eggs
  \item 1 cup milk
\end{ingredients}

\subsection*{Topping}
\begin{itemize}[leftmargin=2em]
  \item Juice of 2 lemons
  \item \nicefrac{1}{2} cup White Sugar
\end{itemize}

\begin{method}
  \item Preheat oven to 200\degree C. Grease or spray a 12-hole muffin tray.
  \item In a large bowl, combine flour, sugar, poppyseeds, and lemon rind.
  \item In a separate bowl, whisk together melted butter, milk, and eggs.
  \item Add wet ingredients to dry and fold until just combined. Do not over-mix.
  \item Spoon mixture into prepared muffin tray. Bake for 10–15 minutes, until golden and springy.
  \item While muffins bake, mix lemon juice with sugar (do not dissolve).
  \item As soon as muffins are removed from the tray, brush lemon sugar glaze over the tops. Cool on a wire rack.
\end{method}

\begin{notes}
The crunchy lemon-sugar glaze gives these muffins a tangy, zesty finish.  
Best enjoyed on the day of baking, but they can be stored in an airtight container for 2–3 days.  
\end{notes}
\end{recipebox}

\section*{Savoury Muffins}
\addcontentsline{toc}{section}{Savoury Muffins}

\begin{recipebox}{Corn Muffins}{12--15}{25 min}
\begin{ingredients}
  \item 1\nicefrac{1}{2} cups self-raising flour
  \item 2 Tbsp grated cheese
  \item Salt and pepper, to taste
  \item 2 eggs
  \item \nicefrac{1}{2} cup milk
  \item 440 g can creamed corn or whole kernel corn
  \item (Optional) Freshly chopped parsley and finely chopped onion
\end{ingredients}

\begin{method}
  \item Preheat oven to 200\degree C.  
  \item Mix flour, cheese, salt, and pepper together in a bowl.  
  \item Beat eggs and milk together in a separate bowl.  
  \item Add the egg mixture and corn to the dry ingredients. Fold together gently, being careful not to over-mix.  
  \item Line a muffin tin with paper cases and spoon in the mixture.  
  \item Bake for 14–16 minutes (medium muffins) or 12–14 minutes (mini muffins), until the tops spring back when gently pressed.  
\end{method}

\begin{notes}
These savoury muffins are perfect for school lunches, morning tea, or served alongside soup.  
For extra flavour, fold in parsley or finely chopped onion before baking.  
\end{notes}
\end{recipebox}

\begin{recipebox}{Spinach and Feta Muffins}{12}{25 min}
\begin{ingredients}
  \item 2 cups self-raising flour
  \item \nicefrac{1}{2} tsp salt
  \item 100 g feta cheese, cut into small cubes
  \item 1 cup spinach leaves, finely chopped
  \item \nicefrac{1}{2} onion, finely chopped
  \item \nicefrac{1}{2} cup grated cheese
  \item 1 egg
  \item 1 cup milk
  \item (Optional) Chilli powder, herbs, pepper to taste
  \item Extra grated cheese or parmesan, for topping
\end{ingredients}

\begin{method}
  \item Preheat oven to 180\degree C. Spray a muffin tray with oil or line with paper cases.
  \item Sift flour and salt into a bowl. Add feta, spinach, onion, and grated cheese.  
  \item Add egg and milk last, folding all ingredients together until just combined.  
  \item Spoon mixture into muffin tray. Sprinkle extra cheese or parmesan on top if desired.  
  \item Bake for 15 minutes, or until golden brown and a toothpick inserted comes out clean.  
\end{method}

\begin{notes}
These savoury muffins are perfect for lunchboxes, picnics, or as a snack.  
For extra flavour, add a pinch of chilli powder, mixed herbs, or black pepper.  
Best enjoyed warm with a spread of butter.  
\end{notes}
\end{recipebox}


